\section{Methods}
\label{sec:methods}

The observation model in the linear Gaussian case is
\begin{equation}
\label{eq:obs}
f_t^{(\theta)}(y_t \mid x_t) = \mathcal{N}(y_t; Hx_t, R).
\end{equation}
An ensemble of size \(N\) is propagated by the state transition. The forecast mean \(\hat{\mu}_{t\mid t-1}\) and covariance \(\hat{\Sigma}_{t\mid t-1}\) are the sample moments, with divisor \(N-1\). The ensemble Kalman likelihood contribution is
\begin{equation}
\label{eq:enkf-lik}
\hat{p}_{\mathrm{enkf},t}^{(\theta)}(y_t \mid y_{0:t-1})
  = \mathcal{N}\bigl(y_t; H\hat{\mu}_{t\mid t-1},\, H\hat{\Sigma}_{t\mid t-1}H^{\top}+R\bigr).
\end{equation}
This is exact for the Gaussian forecast approximation. It is not unbiased for the true predictive density. The product of these contributions, multiplied by the prior, defines \(\hat{\pi}_{\mathrm{enkf},t}\) up to the ensemble randomness that the extended target carries \citep[Section~4.1]{golightly2026}.

\subsection{Nested filter}

The outer recursion follows Algorithm~6 of \citet{golightly2026}. At time \(t\), each parameter particle is updated by one ensemble Kalman step and its weight is multiplied by \eqref{eq:enkf-lik}. If the effective sample size is at least \(\gamma M\), the cloud is left as it is. Otherwise the particles are resampled with replacement, the weights are reset to \(1/M\), and each particle is offered one Metropolis--Hastings proposal. The proposal is a random walk whose covariance is \(\zeta\) times the sample covariance of the other resampled particles. The acceptance ratio uses the prior and the ensemble Kalman likelihood of the series so far. The likelihood at a proposed value is computed by rerunning the filter from the known initial condition. In the pilot, \(\zeta^{2} = 2.56^{2}/d_{\theta}\) and the ensemble size is held fixed.

The printed Algorithm~6 places the effective-sample-size test inside the particle loop. The surrounding text applies one resample-move step when the effective sample size of the updated cloud falls below \(\gamma M\). We follow that reading.

\subsection{Delayed acceptance and the neighbour search}

The screen replaces both ensemble Kalman likelihoods in the acceptance ratio by an inverse-distance average of the \(k\) nearest stored log-likelihoods \citep[Section~4.2]{golightly2026}. If the screen accepts, a second ratio corrects back to the ensemble Kalman likelihood. The product of the two ratios equals the plain Metropolis--Hastings ratio, so the kernel leaves the extended ensemble Kalman target invariant under the conditions of \citet{christen2005}. The distance in the paper is described as some metric. The pilot uses Euclidean distance on \(\log\theta\), after standardising each coordinate by the standard deviation of the stored cloud.

A linear scan and an exact KD-tree that return the same neighbours produce the same surrogate and the same kernel. A different deterministic neighbour set changes the kernel and, if the surrogate is positive wherever the proposal is positive, still leaves the ensemble Kalman target invariant, because the second stage evaluates that target. Approximate neighbours are not used. At the pilot size the question is only whether the tree evaluates a distance to fewer stored points. The tree is rebuilt on every query, so the count is not a claim about wall-clock time.

\subsection{A switch that depends only on time}

While \(t\) is small, the ensemble Kalman likelihood is a short filter and a screen saves little. The pilot rule uses plain Metropolis--Hastings when \(t \le 8\) and delayed acceptance with \(k = 3\) when \(t > 8\), on a series of length 16. The choice is a function of \(t\) only and is made before the proposal. Each kernel preserves \(\hat{\pi}_{\mathrm{enkf},t}\). A mixture of kernels with the same invariant distribution preserves that distribution. The argument does not use the unbiasedness hypothesis in \citet{chopin2012} or the particle-Gibbs kernel in \citet{botha2023}.

The criterion, fixed before any output was seen, is the number of transition draws used inside rejuvenation divided by the number of accepted proposals. That comparison is in the same spirit as the relative cost of the cheap approximation and the stage-two acceptance rate in the delayed-acceptance efficiency analysis of \citet{sherlock2021}. We do not use their diffusion limit.

\subsection{Surrogate incremental weights}

If the incremental weight is the nearest-neighbour prediction of \eqref{eq:enkf-lik}, the weights are aimed at
\begin{equation}
\label{eq:knn-target}
\pi_{\mathrm{knn},t}(\theta) \propto \pi(\theta)\, \exp\bigl(\ell_{\mathrm{knn},t}(\theta)\bigr),
\end{equation}
conditional on the stored cloud that defines \(\ell_{\mathrm{knn},t}\). This is not \(\hat{\pi}_{\mathrm{enkf},t}\) and not the exact posterior. The two errors are separate. Equation~\eqref{eq:enkf-lik} already differs from the true likelihood. The surrogate then differs from \eqref{eq:enkf-lik}.

In the pilot the first \(k\) particles are a reference set and are always updated with the ensemble Kalman filter. Every other particle takes the inverse-distance prediction of the reference increments and does not run a forecast for that weight. After resampling, the reference slots are re-evaluated so that the screen is built from ensemble Kalman likelihoods. A proposal that passes the screen is still evaluated with the ensemble filter. There is no bound on the distance between \(\pi_{\mathrm{knn},t}\) and \(\hat{\pi}_{\mathrm{enkf},t}\). The diagnostic, also fixed in advance, recomputes \eqref{eq:enkf-lik} for eight non-reference particles at the terminal time. A mean, over replicates, of the median absolute log difference above 1 was the threshold for calling the surrogate a poor approximation of the ensemble Kalman increment.

\subsection{Rao-Blackwellised weight}

For a general observation density \(f\), the filter can still be run under a linear Gaussian approximation \(g\). The Rao-Blackwellised increment averages, over the analysis ensemble,
\begin{equation}
\label{eq:rb}
\hat{g}_{\mathrm{enkf},t}^{(\theta)}(y_t)\, \frac{f_t^{(\theta)}(y_t \mid x_t^{j})}{g_t^{(\theta)}(y_t \mid x_t^{j})}.
\end{equation}
The parameter target induced by these increments is denoted \(\tilde{\pi}_{\mathrm{enkf},t}\). If \(f = g\), the ratio equals 1 for every member, and \eqref{eq:rb} equals \eqref{eq:enkf-lik}. The correction does not change the target and the extra density evaluations do not change the estimate. The identity is linear and Gaussian. It is not a general rule for a nonlinear observation model.

Worth its cost, fixed before the runs, means a smaller absolute error of the posterior means of \(\log\theta_i\), per transition draw, than the plain filter that uses the same \(g\). When \(f = g\) the identity says the correction cannot improve the error.
